\section{Giới thiệu tổng quan đề tài}

\subsection{Giới thiệu đề tài}

Trong bối cảnh chuyển đổi số đang diễn ra mạnh mẽ, thư viện truyền thống đang dần chuyển mình để trở thành các trung tâm học liệu trực tuyến hiện đại. Tuy nhiên, thách thức lớn nhất hiện nay không chỉ dừng lại ở việc số hóa các đầu sách mà còn nằm ở cách thức quản lý và khai thác nguồn tài nguyên đó sao cho hiệu quả. Các hệ thống quản lý thư viện hiện có thường tập trung chủ yếu vào các tác vụ quản trị cơ bản, khiến việc tương tác giữa bạn đọc và kho tài liệu vẫn còn mang tính thụ động và rời rạc.

Đối với người quản lý thư viện, việc vận hành thủ công các danh mục tài nguyên khổng lồ từ sách giấy, tạp chí đến các tài nguyên số thường dẫn đến sự chậm trễ trong việc cập nhật trạng thái mượn trả và bảo quản. Sự thiếu hụt các công cụ phân tích tự động khiến công tác quản trị thành viên và theo dõi luồng lưu thông tài liệu trở nên áp lực, đặc biệt là trong việc dự báo nhu cầu đọc để tối ưu hóa bộ sưu tập.

Về phía người đọc, trải nghiệm tìm kiếm tài liệu trên các hệ thống cũ thường gây khó khăn khi chỉ dựa trên các từ khóa khô khan, thiếu sự linh hoạt và cá nhân hóa. Người dùng hiện đại mong muốn một không gian tương tác thông minh hơn, nơi họ có thể nhận được những gợi ý sách phù hợp với sở thích, lịch sử đọc và hành vi cá nhân mà không cần mất quá nhiều thời gian tra cứu thủ công. Đồng thời, nhu cầu được hỗ trợ tức thời thông qua các trợ lý ảo để giải đáp truy vấn về quy định hay vị trí tài liệu đang ngày càng trở nên thiết yếu.

Nhận diện được những yêu cầu cấp thiết đó, đề tài “Phát triển hệ thống quản lý thư viện trực tuyến” được hình thành nhằm nghiên cứu, thiết kế và triển khai một nền tảng công nghệ toàn diện, không chỉ giải quyết triệt để bài toán quản lý tài nguyên và vận hành như danh mục sách, tạp chí, tài nguyên số và quy trình mượn trả mà còn tích hợp Trí tuệ nhân tạo (AI) để nâng cao hiệu quả khai thác bộ sưu tập. Trọng tâm của đề tài là tạo ra một không gian tương tác thông minh dựa trên sự kết hợp giữa mô hình quản trị hiện đại giúp tự động hóa quy trình vận hành và quản lý tập trung tài nguyên, khả năng cá nhân hóa trải nghiệm người dùng thông qua các công cụ tìm kiếm thông minh cùng hệ thống gợi ý tự động dựa trên hành vi đọc thực tế, và sự hỗ trợ tự động hóa từ Chatbot cùng trợ lý ảo nhằm cung cấp dịch vụ giải đáp truy vấn, hướng dẫn sử dụng một cách tức thời cho người dùng mọi lúc, mọi nơi.

Bằng việc ứng dụng các công nghệ phần mềm và thuật toán AI, hệ thống tập trung tối ưu hóa quy trình quản trị cho thủ thư, đồng thời cung cấp cho bạn đọc khả năng tra cứu và tiếp cận tài liệu nhanh chóng, phù hợp với nhu cầu của từng cá nhân.

\subsubsection{Lý do chọn đề tài}

Trong quá trình học tập và sử dụng thư viện tại môi trường đại học, nhóm nhận thấy rằng các hệ thống thư viện hiện tại vẫn còn tồn tại nhiều hạn chế. Việc tìm kiếm tài liệu thường chỉ dựa trên từ khóa đơn giản, dẫn đến kết quả chưa thực sự phù hợp với nhu cầu của người dùng. Bên cạnh đó, hệ thống chưa có khả năng gợi ý tài liệu dựa trên sở thích hoặc lịch sử sử dụng, khiến người đọc phải mất nhiều thời gian để tìm kiếm tài liệu phù hợp.

Ngoài ra, các hoạt động quản lý như theo dõi mượn – trả, cập nhật trạng thái tài liệu hay hỗ trợ người dùng vẫn còn phụ thuộc nhiều vào thao tác thủ công. Điều này không chỉ làm giảm hiệu quả vận hành mà còn ảnh hưởng đến trải nghiệm của người dùng.

Từ những thực tế đó, nhóm quyết định lựa chọn đề tài “Phát triển hệ thống quản lý thư viện trực tuyến” với mong muốn xây dựng một hệ thống hiện đại hơn, có khả năng tự động hóa và hỗ trợ thông minh, góp phần nâng cao hiệu quả quản lý cũng như cải thiện trải nghiệm người dùng.

\subsection{Mục tiêu thực hiện đề tài}

Mục tiêu chính của đề tài là xây dựng một hệ thống quản lý thư viện trực tuyến dưới dạng ứng dụng Web (Web Application) nhằm hỗ trợ hiệu quả cho cả người quản lý thư viện và bạn đọc trong quá trình khai thác và sử dụng tài nguyên học liệu.

Hệ thống được phát triển với mục tiêu cung cấp một nền tảng quản lý tập trung cho các hoạt động vận hành thư viện, bao gồm quản lý danh mục tài nguyên, theo dõi trạng thái mượn – trả tài liệu và quản lý thông tin thành viên. Đồng thời, hệ thống hướng đến việc nâng cao trải nghiệm của người dùng thông qua các công cụ tìm kiếm thông minh và cơ chế gợi ý tài liệu dựa trên sở thích cũng như lịch sử đọc của bạn đọc.

Bên cạnh đó, đề tài cũng hướng đến việc tích hợp các công nghệ trí tuệ nhân tạo nhằm hỗ trợ người dùng trong quá trình tương tác với hệ thống. Các chức năng như chatbot hoặc trợ lý ảo được nghiên cứu và triển khai nhằm cung cấp khả năng giải đáp các câu hỏi thường gặp, hướng dẫn sử dụng hệ thống và hỗ trợ tìm kiếm tài liệu một cách nhanh chóng.

Thông qua việc xây dựng hệ thống này, đề tài hướng tới mục tiêu tạo ra một nền tảng thư viện số thông minh, giúp tối ưu hóa quy trình quản trị cho người quản lý thư viện đồng thời mang lại trải nghiệm truy cập và khai thác tài nguyên thuận tiện, hiệu quả cho bạn đọc.

\subsubsection{Tổng quan hệ thống đề xuất}

Hệ thống quản lý thư viện trực tuyến được xây dựng dưới dạng một ứng dụng web, cho phép người dùng truy cập và sử dụng thông qua trình duyệt mà không cần cài đặt phần mềm bổ sung. Về mặt tổng thể, hệ thống bao gồm các thành phần chính như giao diện người dùng, hệ thống xử lý nghiệp vụ phía máy chủ và cơ sở dữ liệu lưu trữ thông tin tài nguyên cũng như người dùng.

Bên cạnh các thành phần cơ bản, hệ thống còn tích hợp các mô-đun trí tuệ nhân tạo nhằm hỗ trợ các chức năng nâng cao như gợi ý tài liệu, tìm kiếm thông minh và hỗ trợ người dùng thông qua chatbot. Các mô-đun này hoạt động dựa trên dữ liệu thu thập từ hành vi người dùng và thông tin tài nguyên trong hệ thống, từ đó đưa ra các kết quả phù hợp và có tính cá nhân hóa.

\subsection{Lợi ích của đề tài mang lại}

Đề tài hướng đến việc mang lại nhiều lợi ích thiết thực cho hoạt động quản lý và khai thác tài nguyên thư viện.

\subsubsection{Xây dựng hệ thống quản lý thư viện trực tuyến hiện đại}

Đề tài tập trung xây dựng một hệ thống quản lý thư viện dưới dạng ứng dụng web, cho phép quản lý và truy cập tài nguyên thông qua môi trường trực tuyến. Hệ thống tích hợp các công nghệ xử lý dữ liệu và trí tuệ nhân tạo nhằm nâng cao khả năng quản lý cũng như hỗ trợ người dùng trong quá trình tìm kiếm và khai thác tài liệu.

\subsubsection{Hỗ trợ hiệu quả cho công tác quản lý thư viện}

Hệ thống cho phép quản lý danh mục tài nguyên đa dạng như sách, tạp chí và tài nguyên số. Các chức năng theo dõi trạng thái mượn – trả – bảo quản tài liệu, quản lý thông tin thành viên và thống kê dữ liệu sử dụng giúp người quản lý thư viện vận hành hệ thống một cách hiệu quả hơn. Các dữ liệu thống kê từ hệ thống cũng hỗ trợ việc đưa ra các quyết định liên quan đến việc bổ sung hoặc điều chỉnh bộ sưu tập tài liệu.

\subsubsection{Tối ưu hóa việc khai thác tài liệu của bạn đọc}

Các chức năng tìm kiếm thông minh, gợi ý tài liệu tự động và hệ thống trợ lý ảo giúp người dùng có thể nhanh chóng tìm thấy các tài liệu phù hợp với nhu cầu học tập và nghiên cứu. Hệ thống cũng hỗ trợ cá nhân hóa trải nghiệm người dùng dựa trên lịch sử đọc và hành vi tìm kiếm, từ đó giúp nâng cao hiệu quả khai thác tài nguyên học liệu.

\subsubsection{Nâng cao trải nghiệm tương tác giữa người dùng và hệ thống thư viện}

Việc tích hợp chatbot hoặc trợ lý ảo giúp người dùng có thể nhận được sự hỗ trợ nhanh chóng khi sử dụng hệ thống. Các truy vấn liên quan đến quy định mượn sách, tìm kiếm tài liệu hoặc hướng dẫn sử dụng hệ thống có thể được giải đáp một cách tức thời, góp phần cải thiện trải nghiệm tổng thể của người dùng khi tương tác với thư viện.

\subsection{Phạm vi và giới hạn của đề tài}

Phạm vi của đề tài tập trung vào việc nghiên cứu và xây dựng hệ thống quản lý thư viện trực tuyến dưới dạng ứng dụng Web (Web Application). Các đối tượng tương tác trực tiếp với hệ thống bao gồm: Quản trị viên (Admin), Thủ thư (Librarian) và Bạn đọc (Reader).

\subsubsection{Các nhóm người dùng}

Bạn đọc (Reader): đối tượng sử dụng chính của hệ thống, có nhu cầu tìm kiếm, tra cứu và mượn tài liệu. Người dùng cũng có thể nhận các gợi ý tài liệu phù hợp, quản lý lịch sử mượn sách và tương tác với hệ thống thông qua các chức năng tìm kiếm hoặc chatbot.

Thủ thư (Librarian): chịu trách nhiệm quản lý danh mục tài liệu, xử lý các hoạt động mượn – trả và cập nhật trạng thái tài nguyên. Ngoài ra, thủ thư cũng có thể sử dụng các công cụ hỗ trợ để nhập dữ liệu và kiểm soát thông tin trong hệ thống.

Quản trị viên (Admin): có vai trò quản lý toàn bộ hệ thống, bao gồm cấu hình các thông số, quản lý tài khoản người dùng và giám sát hoạt động chung của hệ thống.

Hệ thống cho phép quản lý các loại tài nguyên thư viện phổ biến như sách in, tạp chí vật lý và các tài nguyên số dưới dạng tệp tin điện tử như E-book hoặc tài liệu PDF. Các chức năng chính của hệ thống bao gồm quản lý danh mục tài nguyên, quản lý thông tin thành viên, theo dõi trạng thái mượn – trả tài liệu, hỗ trợ tìm kiếm tài liệu và cung cấp các gợi ý tài liệu phù hợp với nhu cầu của người dùng.

Trong phạm vi đề tài, hệ thống được thiết kế hướng đến các thư viện quy mô vừa và nhỏ như thư viện khoa, thư viện trường học hoặc các trung tâm học liệu. Đối tượng người dùng chính bao gồm sinh viên, giảng viên và các nhà nghiên cứu có nhu cầu truy cập và khai thác tài liệu học thuật.

Các tính năng AI được triển khai ở mức cơ bản, chi tiết được trình bày ở phần phạm vi ứng dụng trí tuệ nhân tạo bên dưới.

\subsubsection{Phạm vi ứng dụng trí tuệ nhân tạo}

Trong phạm vi của đề tài, các chức năng liên quan đến trí tuệ nhân tạo được triển khai ở mức cơ bản nhằm minh họa khả năng tích hợp vào hệ thống quản lý thư viện. Cụ thể, hệ thống sử dụng các mô hình và dịch vụ có sẵn để xây dựng các chức năng như gợi ý tài liệu dựa trên hành vi người dùng, tìm kiếm bằng ngôn ngữ tự nhiên và chatbot hỗ trợ trả lời các câu hỏi phổ biến.

Do giới hạn về thời gian và nguồn lực, đề tài không tập trung vào việc xây dựng hoặc huấn luyện các mô hình AI từ đầu, mà chủ yếu tận dụng các công cụ và API hiện có để triển khai và tích hợp vào hệ thống. Điều này giúp đảm bảo tính khả thi của đề tài trong khi vẫn thể hiện được hướng tiếp cận hiện đại.

\subsubsection{Giới hạn của đề tài}

Mặc dù hệ thống được thiết kế với nhiều chức năng, đề tài vẫn tồn tại một số giới hạn nhất định. Hệ thống hiện chỉ được phát triển dưới dạng ứng dụng web và chưa hỗ trợ ứng dụng trên nền tảng di động. Các chức năng trí tuệ nhân tạo được triển khai ở mức cơ bản và chưa đạt đến mức tối ưu như trong các hệ thống thương mại lớn.

Ngoài ra, hệ thống chưa xử lý các bài toán quy mô lớn như dữ liệu hàng triệu bản ghi hoặc tích hợp sâu với các hệ thống thư viện bên ngoài. Các tính năng nâng cao sẽ được xem xét phát triển trong các nghiên cứu và phiên bản tiếp theo.

\subsection{Ý nghĩa của đề tài}

\subsubsection{Ý nghĩa thực tiễn}

Hệ thống quản lý thư viện trực tuyến được phát triển trong đề tài góp phần nâng cao hiệu quả vận hành của các thư viện trong môi trường số. Việc quản lý tập trung các loại tài nguyên và thông tin người dùng giúp giảm thiểu các thao tác thủ công, đồng thời đảm bảo dữ liệu luôn được cập nhật và lưu trữ một cách chính xác.

Bên cạnh đó, hệ thống còn giúp cải thiện đáng kể trải nghiệm của bạn đọc thông qua các công cụ tìm kiếm và gợi ý tài liệu thông minh. Người dùng có thể nhanh chóng tìm thấy những tài liệu phù hợp với nhu cầu học tập và nghiên cứu mà không cần thực hiện quá nhiều thao tác tra cứu.

Ngoài ra, việc tích hợp chatbot hoặc trợ lý ảo trong hệ thống cũng giúp cung cấp dịch vụ hỗ trợ tức thời cho người dùng, đồng thời giảm tải khối lượng công việc cho nhân viên thư viện trong việc giải đáp các câu hỏi phổ biến.

\subsubsection{Ý nghĩa khoa học}

Về mặt nghiên cứu, đề tài góp phần làm rõ cách thức xây dựng và triển khai một hệ thống quản lý thư viện dựa trên các công nghệ phần mềm hiện đại. Quá trình thiết kế hệ thống bao gồm việc xây dựng kiến trúc phần mềm, thiết kế cơ sở dữ liệu và triển khai các chức năng quản lý tài nguyên và người dùng.

Đề tài cũng nghiên cứu việc tích hợp các mô hình trí tuệ nhân tạo trong hệ thống thông tin nhằm hỗ trợ cá nhân hóa trải nghiệm người dùng. Việc ứng dụng các thuật toán gợi ý tài liệu và chatbot hỗ trợ người dùng giúp minh họa tiềm năng của trí tuệ nhân tạo trong việc nâng cao hiệu quả khai thác tài nguyên tri thức.

Thông qua việc kết hợp giữa các công nghệ phát triển ứng dụng web, quản lý dữ liệu và trí tuệ nhân tạo, đề tài góp phần đề xuất một mô hình hệ thống quản lý thư viện thông minh có thể được tham khảo và mở rộng trong các nghiên cứu và ứng dụng thực tế trong tương lai.

Trong các chương tiếp theo, đề tài sẽ tập trung vào việc phân tích yêu cầu hệ thống, thiết kế kiến trúc tổng thể và hiện thực các chức năng cốt lõi nhằm xây dựng một hệ thống quản lý thư viện hoàn chỉnh.